% ============ §21 Complete intersection rings (原书 pp.169-173) ============
\section{Complete intersection rings}\label{sec:21}\index{complete intersection rings, section 21}\index{c.i. ring, section 21}

Let $(A, m, k)$ be a Noetherian local ring; we choose a minimal basis
$x_1, \ldots, x_n$ of $m$, where $n = \embdim A$ is the embedding dimension (see
\S\ref{sec:14}). Set $E_\bullet = K_\bullet(x_1, \ldots, x_n)$ for the Koszul complex. The complex $E_\bullet$ is
determined by $A$ up to isomorphism. Indeed, if $x_1', \ldots, x_n'$ is another
minimal basis of $m$ then by Theorem~\ref{thm:2.3}, there is an invertible $n \times n$
matrix $(a_{ij})$ over $A$ such that $x_i = \sum a_{ij}x_j'$. It is proved in Appendix C that
$K_\bullet(x_1, \ldots, x_n)$ can be thought of as the exterior algebra $\bigwedge(Ae_1 + \cdots + Ae_n)$ with
differential defined by $d(e_i) = x_i$. Similarly,
$K_\bullet(x_1', \ldots, x_n') = \bigwedge(Ae_1' + \cdots + Ae_n')$ with $d(e_i') = x_i'$.
Now $f(e_j') = \sum a_{ij}e_i$ defines an isomorphism from the free $A$-module
$Ae_1' + \cdots + Ae_n'$ to $Ae_1 + \cdots + Ae_n$, which extends to an isomorphism $f$
of the exterior algebra; $f$ commutes with the differential $d$, since for a
generator $e_i'$ of $\bigwedge(Ae_1' + \cdots + Ae_n')$ we have $df(e_i') = \sum a_{ij}x_j = x_i = fd(e_i')$.
Therefore $f\colon K_\bullet(x_1', \ldots, x_n') \to K_\bullet(x_1, \ldots, x_n)$ is an isomorphism of complexes.

Since $mH_p(E_\bullet) = 0$ by Theorem~\ref{thm:16.4}, $H_p(E_\bullet)$ is a vector space over
$k$, and we set
\[
  \varepsilon_p = \dim_k H_p(E_\bullet) \quad \text{for } p = 0, 1, 2, \ldots;
\]
then these are invariants of the local ring $A$. In view of $H_0(E_\bullet) = A/(x) =
A/m = k$, we have $\varepsilon_0 = 1$. In this section we are concerned with $\varepsilon_1$. If $A$ is
regular then $x_1, \ldots, x_n$ is an $A$-sequence, so that $\varepsilon_1 = \cdots = \varepsilon_n = 0$, and
conversely by Theorem~\ref{thm:16.5}, $\varepsilon_n = 0$ implies that $A$ is regular.

Let us consider the case when $A$ can be expressed as a quotient of a
regular local ring $R$; let $A = R/a$, and write $n$ for the maximal ideal of $R$.
If $a \nsubseteq n^2$, we can take $x \in a - n^2$; then $R' = R/xR$ is again a regular local
ring, and $A = R'/a'$, so that we can write $A$ as a quotient of a ring of
dimension smaller than $R$. In this way we see that there exists an
expression $A = R/a$ of $A$ as a quotient of a regular local ring $(R, n)$ with
$a \subset n^2$. Then we have $m = n/a$ and $m/m^2 = n/(a + n^2) = n/n^2$, so that
$\dim R = \embdim A$. Conversely, equality here implies that $a \subset n^2$.

Let $(R, n)$ be a regular local ring and $A = R/a$ with $a \subset n^2$; choose a
regular system of parameters (that is a minimal basis of $n$) $t_1, \ldots, t_n$.
Then the images $x_i$ of $t_i$ in $A$ form a minimal basis $x_1, \ldots, x_n$ of $m$. Let
\[
  K_\bullet(t_1, \ldots, t_n)\colon 0 \to L_n \to L_{n-1} \to \cdots \to L_1 \to L_0 \to 0
\]
be the Koszul complex of $R$ and $t$. By Theorem~\ref{thm:16.5}, we know that this
becomes exact on adding $\cdots + L_0 \to k \to 0$ to the right-hand end, so that
$K_\bullet(t_1, \ldots, t_n)$ is a projective resolution of $k$ as an $R$-module. Taking the tensor
product with $A = R/a$, we get the complex $E_\bullet = K_\bullet(x_1, \ldots, x_n)$ of $A$-modules.
Thus we have
\[
  H_p(E_\bullet) = H_p(K_\bullet(t_1, \ldots, t_n) \otimes_R A) = \Tor_p^R(k, A) \quad \text{for all } p \ge 0.
\]
However, from the exact sequence of $R$-modules $0 \to a \to R \to A \to 0$ we
get the long exact sequence
\[
  0 = \Tor_1^R(k, R) \to \Tor_1^R(k, A) \to k \otimes_R a \to k \otimes_R R
      \to k \otimes_R A \to 0;
\]
at the right-hand end we have $k \otimes_R R \cong k \otimes_R A = k$, so that
\[
  \Tor_1^R(k, A) \cong k \otimes_R a = a/na.
\]
Quite generally, we write $\mu(M)$ for the minimum number of generators of
an $R$-module $M$. Then we see that
\[
  \mu(a) = \dim_k H_1(E_\bullet) = \varepsilon_1(A).
\]

\begin{theorem}\label{thm:21.1}
Let $(A, m, k)$ be a Noetherian local ring, and $\hat{A}$ its
completion.
\begin{enumerate}[(i)]
  \item $\varepsilon_p(A) = \varepsilon_p(\hat{A})$ for all $p \ge 0$.
  \item $\varepsilon_1(A) \ge \embdim A - \dim A$.
  \item If $R$ is a regular local ring, $a$ an ideal of $R$ and $A \cong R/a$, then
        $\mu(a) = \dim R - \embdim A + \varepsilon_1(A)$.
\end{enumerate}
\end{theorem}

\begin{proof}
(i) is clear from the fact that a minimal basis of $m$ is a minimal
basis of $m\hat{A}$, so that applying $\otimes_A \hat{A}$ to the complex $E_\bullet$ made from $A$
gives that made from $\hat{A}$. Then since $\hat{A}$ is $A$-flat, $H_p(E_\bullet) \otimes \hat{A} = H_p(E_\bullet \otimes \hat{A})$,
and $mH_p(E_\bullet) = 0$ gives $H_p(E_\bullet) \otimes \hat{A} = H_p(E_\bullet)$.

(ii) If $A$ is a quotient of a regular local ring, then as we have seen above,
there exists a regular ring $(R, n)$ such that $A = R/a$ with $a \subset n^2$, so that
$\varepsilon_1(A) = \mu(a) \ge \htop a = \dim R - \dim A = \embdim A - \dim A$, where the
equality for $\htop a$ comes from Theorem~\ref{thm:17.4}, (i). Now $A$ itself is not necessarily
a quotient of a regular local ring, but we will prove later (see \S29) that $\hat{A}$
always is, and we admit this in the section. Having said this, the two sides
of (ii) are unaltered on replacing $A$ by $\hat{A}$, and the inequality holds for $\hat{A}$.

(iii) Set $n = \rad(R)$. If $a \subset n^2$ then, as we have seen above, $\mu(a) = \varepsilon_1(A)$,
and $\dim R = \embdim A$, so that we are done. If $a \nsubseteq n^2$, take $x \in a - n^2$;
we pass to $R/xR$ and $a/xR$, each of $\dim R$ and $\mu(a)$ decreases by 1, so
induction completes the proof.
\end{proof}

\begin{definition*}
A Noetherian local ring $A$ is a \emph{complete intersection ring}
(abbreviated to \emph{c.i. ring}) if $\varepsilon_1(A) = \embdim A - \dim A$.
\end{definition*}

\begin{theorem}\label{thm:21.2}
Let $A$ be a Noetherian local ring.
\begin{enumerate}[(i)]
  \item $A$ is c.i. $\iff$ $\hat{A}$ is c.i.
  \item Let $A$ be a c.i. ring and $R$ a regular local ring such that $A = R/a$;
        then $a$ is generated by an $R$-sequence. Conversely, if $a$ is an ideal generated
        by an $R$-sequence then $R/a$ is a c.i. ring.
  \item A necessary and sufficient condition for $A$ to be a c.i. ring is that
        the completion $\hat{A}$ should be a quotient of a complete regular local ring
        by an ideal generated by an $R$-sequence.
\end{enumerate}
\end{theorem}

\begin{proof}
(i) is obvious.

(ii) By Theorem~\hyperref[thm:21.1]{1}, (iii), $\mu(a) = \dim R - \embdim A + \varepsilon_1(A)$, and by
Theorem~\ref{thm:17.4}, (i), $\htop a = \dim R - \dim A$, so that $A$ is a c.i. ring is equivalent
to $\htop a = \mu(a)$. But by Theorem~\ref{thm:17.4}, (iii), this is equivalent to $a$ being
generated by an $R$-sequence.

(iii) The sufficiency is clear from (i) and (ii). Necessity follows from the
fact that $\hat{A}$ is a quotient of a complete regular local ring (see \S29), together
with (i) and (ii).
\end{proof}

\begin{theorem}\label{thm:21.3}
A c.i. ring is Gorenstein.
\end{theorem}

\begin{proof}
If $A$ is c.i. then so is $\hat{A}$, and if $\hat{A}$ is Gorenstein then so is $A$, so that we
can assume that $A$ is complete. Then we can write $A = R/a$, where $R$ is a
regular local ring and $a$ is an ideal generated by a regular sequence. Since $R$
is Gorenstein, $A$ is also by Ex.~\ref{ex:18.1-ex}.
\end{proof}

Thus we have the following chain of implications for Noetherian local
rings:
\[
  \text{regular} \implies \text{c.i.} \implies \text{Gorenstein} \implies \text{CM}.
\]
Let $A$ be a c.i. ring, and $p$ a prime ideal of $A$. If $A$ is of the form
$R/(x_1, \ldots, x_r)$, where $R$ is regular and $x_1, \ldots, x_r$ is an $R$-sequence, then
since $A_p$ can be written $A_p = R_p/(x_1, \ldots, x_r)$, where $R_p$ is regular and
$x_1, \ldots, x_r$ is an $R_p$-sequence, it follows that $A_p$ is again a c.i. ring. The
question of deciding whether $A_p$ is still a c.i. ring even if $A$ is not a quotient
of a regular local ring remained unsolved for some time, but was answered
affirmatively by Avramov~\cite{Av1}, making use of Andr\'e's homology theory
[1, 2]. % 注: 原书此处引 [1, 2]; André 在文献表中仅一条论文条目, 故保持纯文本
This theory defines homology and cohomology groups
$H_n(A, B, M)$ and $H^n(A, B, M)$ for $n \ge 0$ associated with a ring $A$, an $A$-algebra $B$
and a $B$-module $M$. The definition is complicated, but in any case these are
$B$-modules having various nice functorial properties. If $A$ is a
Noetherian local ring with residue field $k$ then
\[
  A \text{ is regular} \iff H_2(A, k, k) = 0,
\]
and
\[
  A \text{ is c.i.} \iff H_3(A, k, k) = 0;
\]
for $n \ge 3$ the statements $H_3(A, k, k) = 0$ and $H_n(A, k, k) = 0$ are equivalent.
Thus Andr\'e homology is particularly relevant to the study of regular and
c.i. rings.

% ============ 习题 ============
\section*{Exercises to \S\ref{sec:21}}
\emph{Prove the following propositions.}

\begin{exercise}\label{ex:21.1-ex}
Let $R$ be a regular ring, $I$ an ideal of $R$, and let $A = R/I$; then the subset
$\{p \in \Spec A \mid A_p \text{ is c.i.}\}$ is open in $\Spec A$ (use Theorem~\ref{thm:19.9}).
\end{exercise}

\begin{exercise}\label{ex:21.2-ex}
Let $A$ be a Noetherian local ring with $\embdim A = \dim A + 1$; if $A$ is CM
then it is c.i.
\end{exercise}

\begin{exercise}\label{ex:21.3-ex}
Let $k$ be a field, and set $A = k[X, Y, Z]/(X^2 - Y^2, Y^2 - Z^2, XY, YZ, ZX)$;
then $A$ is Gorenstein but not c.i.
\end{exercise}
